% Research in progress.
\subsection{The cocycle, its primitive, and the manuscript's splittings}

\begin{statement}{D.1}{The conditional manuscript dictionary}{Proved algebraically; manuscript identification conditional}
Let $F$ be the profinite tripod group with peripheral generators $l_0,l_1$,
let $\rho$ interchange their labels, and let a normalized continuous
automorphism $u$ satisfy
\[
u(l_0)=l_0^p,\qquad u\rho u^{-1}=\operatorname{int}(g)\rho,
\qquad \rho(l_0)=l_1.
\]
Then, in the convention $u(l_1)=f^{-1}l_1^\lambda f$,
we have $\lambda=p$ and $f=g^{-1}$. If a lift of $\rho$ in a semidirect
extension is moved by a frame $A$, then
\[
u(\rho)=A\rho A^{-1},\qquad g=A\rho(A)^{-1}.
\tag{D1.1}
\]
The cocycle $g$ is therefore not itself the frame. Replacing $u$ by
$\operatorname{int}(l_0^a)u$ changes $g$ to
$l_0^a g l_1^{-a}$, for $a\in\wh\Z$.

Conditional on the readings in the supplied note
\emph{EXPERT\_ON\_INTAKE\_11}, these are precisely the equations of
cote 145, scans 7 and 11, and the frames of LM7 D.25 and LM7\_PUSH E.65
move the canonical isotropy splittings described in scans 22 and 28--29,
at the equianharmonic point $P$ and the harmonic points $2,-1,1/2$.
The original scans of cote 145 were not supplied or independently read.
The identification of its $(\gamma,\chi)$ with $(f,\lambda)$ is therefore
conditional on that manuscript reading and the displayed gauge convention.
No identification of the separate symbol $\beta$ on scan 11 is established.
\end{statement}
\begin{proof}
Applying $u\rho=\operatorname{int}(g)\rho u$ to $l_0$ gives
$u(l_1)=g l_1^p g^{-1}$, which proves the convention comparison.
Inside $F\rtimes\langle\rho\rangle$, the product
$A\rho A^{-1}$ equals $A\rho(A)^{-1}\rho$; this proves (D1.1).
Substitution into the same equation after conjugating $u$ by $l_0^a$
gives the stated gauge action.

For an isotropy subgroup $H$ with a chosen splitting $j:H\to E$, a
frame means an element $A$ such that the transported splitting is
$\operatorname{int}(A)\circ j$ (with the specified automorphism of $H$
when one is present). Two frames differ on the right by an element of
$C_E(j(H))$ in the allowed geometric subgroup. This types the ambiguity:
a frame is a conjugator, and the resulting cocycle is its multiplicative
coboundary under the isotropy action. The same algebra applies to the
order-three and order-two splittings separately. The note explicitly
records those canonical isotropy splittings at the named manuscript
locations, but no algebraic calculation verifies a missing palaeographic
reading. We consequently retain that dependence in the statement itself.
In particular, an exponent appearing in $g$ cannot be identified with an
abelianized coordinate of a primitive $A$ without another equation.
\end{proof}

\subsection{The questions of Sections 51.15 and 51.16}

The quotations below are from the supplied AI-made transcription of
cote 140-4, not from newly inspected manuscript images. We retain its
section and scan numbers. Its folio markers place scans 171, 173, 175
and 177 at folios 733, 734, 735 and 736 respectively; hence the intervening
scans are their versos. In particular, scan 178 is folio 736v on this
numbering. The prompt's shorthand ``733v--735, scans 172--178'' does not
cover the last two marked folios. We do not silently substitute that
shorthand for the transcription's actual markers.
The transcription's (357) labels the Tate module of $\mathbb G_m$
by the older twist convention $(-1)$, retained in (358)--(365).
Our modern convention calls the same module $\widehat\Z(1)$.
D.3--D.4 translate that convention; the underlying homomorphisms on
the cyclotomic kernel are unchanged.

\begin{statement}{D.2}{The answer beside each original question}{Proved as a provenance comparison}
The following passages state the original questions answered by the
specified results of Parts A and B. Each answer has exactly the field,
quotient and local hypotheses of those results; a quotient calculation
does not answer the full extension question.
\end{statement}
\begin{proof}
In Section 51.16, folio 733, scan 171, equation (366), the curve is the
\emph{plus} curve $x_0^n+x_1^n+x_\infty^n=0$ with all coordinates nonzero.
The transcription says that it ``n'a pas de pt rat $/\mathbb Q$''.
For odd $n$ the substitution $z=-x_\infty$ identifies it with the
classical curve. For even $n$ the plus curve has no real point at all,
whereas the classical curve has four rational cusps. Statements A.1--A.16,
A.20--A.21 and A.31 give the section answer, keeping this distinction:
the odd full Selmer section set is the union of the three rational cusp
packets; the classical even answer is the union of four packets; the plus
even answer is empty. The assertions for all sections use the named local
Theorem 1.1. The local calculations A.3--A.6 and A.28--A.30 identify the
arithmetic inputs needed for this answer rather than adding a new reading
to the manuscript.

In Section 51.16, folio 733v, scan 172, equation (371), the text distinguishes
the splittings that are ``distincts des `scindages \`a l'infini' ''.
On folio 734, scan 173, it asks whether one can prove
``l'inexistence de scindages de (371) autres que ceux `provenant des pts
\`a l'infini' ''. The same page proposes ``des `approximations' de plus
en plus fines de $\Pi_1(U_n)$ et de la structure d'extension (371)''.
The repaired full Selmer theorem in A.1--A.16 answers the first question
for the Fermat curves in question, and A.20--A.21 give the exact additional
local input for all full sections. The twist results A.23--A.27 and A.32--A.43
extend that answer only in their stated rank-zero classes. They retain
the infinitely many classes in each cusp packet, a datum absent from a
mere list of cusp locations.

The quotient question is explicit in Section 51.16, folio 735, scan 175,
equation (373), and folios 735v--736, scans 176--177, equations (377)--(379).
On scan 177 the text asks to ``d\'eterminer toutes ces classes de
$\widehat M[n]$-conjugaison possibles de tels scindages'' and then asks
which can correspond to a rational point. Statements B.1--B.22 compare
the full, relative prosolvable, maximal prosolvable, derived, pro-$3$
nilpotent and metabelian problems with this question. The subgroup and
coefficient distinctions in those statements are essential: the cubic
abelian Selmer result does not identify all abelian sections, and the
infinitely ramified class-three construction is not a full profinite
section. The unresolved Ross value, intrinsic action, general metabelian
classification and colouring remain unanswered parts of this original
request for finer approximations. B.15--B.17 give the requested explicit
phantom coordinates and the fibre comparison. B.19--B.22 separate the
known Hodge regulator from the unresolved arithmetic action and correct
the scope of its proposed integral test.

Finally Section 51.16, folio 736v, scan 178, says
``l'approximation sur laquelle nous travaillons est beaucoup trop
grossi\`ere'' and concludes that ``il est n\'ecessaire de passer \`a des
approximations plus fines''. B.2's relative prosolvable quotient gives one
precise refinement that carries the Fermat theorem. B.7 and B.15--B.17
show why the existence of a section in a nilpotent
or metabelian approximation alone does not give a section of (371).
No extrapolation along a noncofinal tower is used in this comparison.

For Section 51.15, folio 731v, scan 168, questions (358)(a)--(b) ask whether
the Kummer maps attached to distinct primes are surjective or have open
image, separately and jointly. Folio 732, scan 169, decomposes them into
coefficient primes; folio 732v, scan 170, asks the stronger local-inertia
question (365). A.29 and A.32 use exactly the completed Kummer parameters
whose components this passage distinguishes. D.3 below answers the joint
global image question. It does not claim to compute every local-inertia
image in (365). This is also the historical source of the coefficient and
completion distinctions retained throughout Part B.
Statement D.4 below subsequently determines that local-inertia image
itself, including its exceptional factor at the coefficient prime two.

There is also a mathematical inconsistency on scan 178: the displayed
affine equation there has right side $1$, whereas equation (366) defines
the plus curve by the homogeneous sum equal to $0$. We use (366) and
the actual coordinate-ratio embedding in the torus, as proved in A.17;
we do not use the inconsistent right side to identify a cover or its
arithmetic fundamental group.
\end{proof}

\begin{statement}{D.3}{The joint global Kummer image in Section 51.15}{Proved}
Let $p_1,\ldots,p_m$ be distinct positive rational primes and put
$G_{\Q}^{+}=\Gal(\overline\Q/\Q^{\mathrm{ab}})=\ker\chi$.
Fix compatible roots and compatible roots of unity. The joint Kummer map
\[
G_{\Q}^{+}\longrightarrow\wh\Z(1)^m
\]
has image $2\wh\Z(1)^m$. Thus it has index $2^m$, is surjective at every
odd coefficient prime, and has image $2\Z_2(1)^m$ at $2$. In particular,
the surjectivity alternative in Section 51.15, scan 168, is false, while
its open-image alternative holds with this exact index. This concerns
the full cyclotomic kernel, not the smaller local inertia group of (365).
\end{statement}
\begin{proof}
All square roots of rational numbers lie in $\Q^{\mathrm{ab}}$, so every
Kummer coordinate is even. To prove that these are the only restrictions,
work at a finite level $n$ over the field $K=\Q^{\mathrm{ab}}$, which
contains $\mu_n$. Kummer duality says that the annihilator of the image
in $(\Z/n)^m$ consists of the exponent vectors $(e_i)$ for which
$a=\prod_i p_i^{e_i}$ has an $n$th root in $K$.

Choose the positive real root $\alpha=a^{1/n}$. If $\alpha\in K$, complex
conjugation commutes on $K$ with every $\sigma\in G_\Q$. Writing
$\sigma(\alpha)=\zeta\alpha$ gives both
$c\sigma(\alpha)=\zeta^{-1}\alpha$ and
$\sigma c(\alpha)=\zeta\alpha$. Therefore $\zeta^2=1$ for every
$\sigma$, so $\alpha^2\in\Q$. Unique factorization in $\Q^\times$ now
gives $n\mid2e_i$ for every $i$. Conversely, those divisibilities make
$\alpha^2$ rational, so $\alpha$ is quadratic over $\Q$ and belongs to
$K$. The annihilator is therefore exactly
\[
\{(e_i)\in(\Z/n)^m:2e_i=0\ \text{for every }i\}.
\]
Its orthogonal subgroup is $2(\Z/n)^m$. The image of the compact group
$G_\Q^+$ is closed in $\wh\Z^m$, and its reductions have now been
determined for every $n$. It must equal $2\wh\Z^m$. This proof also
shows independence of the selected primes; checking their images one at
a time would not suffice for the joint assertion.
\end{proof}

\IfFileExists{tables/manuscript_concordance.tex}{\input{tables/manuscript_concordance.tex}}{}

\begin{statement}{D.4}{The exact local-inertia Kummer image}{Proved}
Fix a rational prime $\ell$, a decomposition group at $\ell$, its
inertia subgroup $I_\ell$, and
$I_\ell^+=\ker(\chi_\ell:I_\ell\to\Z_\ell^\times)$.
For a positive rational prime $p$, the image of the continuous Kummer
homomorphism $\kappa_\ell(p)|I_\ell^+$ is
\[
\begin{cases}
 \Z_\ell(1),&\ell\ne2, p=\ell,\\
 \ell^{v_\ell(p^{\ell-1}-1)-1}\Z_\ell(1),&\ell\ne2, p\ne\ell,\\
 2\Z_2(1),&\ell=p=2,\\
 2^{v_2(p^2-1)-2}\Z_2(1),&\ell=2, p\ne2.
\end{cases}
\tag{D4.1}
\]
Thus the open-image question in Section 51.15, folio 732v, scan 170,
equation (365), has an affirmative answer for each $p,\ell$.
Surjectivity is false in general, and always fails at $\ell=2$.

More generally, for $a_1,\ldots,a_m\in\Q_\ell^\times$, write
\[
 a_i=\ell^{v_i}\omega_i(1+\ell)^{d_i}\quad(\ell\ne2),
 \qquad a_i=2^{v_i}\epsilon_i5^{d_i}\quad(\ell=2),
\]
where $d_i\in\Z_\ell$, $\omega_i\in\mu_{\ell-1}$ and
$\epsilon_i\in\{1,-1\}$. Put $v=(v_i)$ and $d=(d_i)$.
After a compatible identification of the Tate twists, the joint image is
\[
 \Z_\ell v+\Z_\ell d\quad(\ell\ne2),
 \qquad 2\Z_2v+2\Z_2d\quad(\ell=2).
\tag{D4.2}
\]
In particular its rank is at most two, and it is not open in
$\Z_\ell(1)^m$ when $m>2$. The joint local assertion is therefore
strictly different from the joint global assertion D.3.
\end{statement}
\begin{proof}
The local Kronecker--Weber theorem identifies
\[
 \Q_\ell^{\mathrm{nr}}(\mu_{\ell^\infty})
   =\Q_\ell^{\mathrm{ab}},
 \qquad I_\ell^+=G_{\Q_\ell^{\mathrm{ab}}};
\]
see \citet[I, Theorem 4.8 and Corollary 4.12]{BMilneCFT2020},
with the multiplicative Lubin--Tate group for $\Q_\ell$.
We first calculate which radicals of an element $a\in\Q_\ell^\times$
can lie in this abelian extension.

If $a^{1/\ell^n}\in\Q_\ell^{\mathrm{ab}}$, its finite Kummer cocycle
$c$ factors through the abelian group $G_{\Q_\ell}^{\mathrm{ab}}$.
The cocycle identity and commutativity give, for every $g,h$,
\[
 (\chi_\ell(g)-1)c(h)=(\chi_\ell(h)-1)c(g)
 \quad\text{in }\Z/\ell^n(1).
\tag{D4.3}
\]
For odd $\ell$, choose $g$ with $\chi_\ell(g)=2$.
Then $c(h)=(\chi_\ell(h)-1)c(g)$ is a coboundary. Kummer theory
implies $a\in(\Q_\ell^\times)^{\ell^n}$; the converse is immediate.
For $\ell=2$, choose $g$ with $\chi_2(g)=-1$. Equation (D4.3)
shows that $2c$ is a coboundary, hence
$a^2\in(\Q_2^\times)^{2^n}$, or equivalently
\[
 a\in\{1,-1\}(\Q_2^\times)^{2^{n-1}}.
\tag{D4.4}
\]
Conversely, if $a=\epsilon b^{2^{n-1}}$, a $2^n$th root is the
product of a square root of $b$ and a $2^n$th root of $\epsilon$.
Both lie in $\Q_2^{\mathrm{ab}}$, so (D4.4) is exact.

At level $\ell^n$, the annihilator of the joint Kummer image consists
of $e\in(\Z/\ell^n)^m$ such that
$\prod_i a_i^{e_i}$ has an $\ell^n$th root in
$\Q_\ell^{\mathrm{ab}}$. For odd $\ell$, the preceding calculation
says precisely $v\cdot e=d\cdot e=0$ modulo $\ell^n$.
For $\ell=2$, it says $2v\cdot e=2d\cdot e=0$ modulo $2^n$;
the signs contribute no condition. Finite-module duality identifies
the orthogonal subgroup with the span of $v,d$, or of $2v,2d$,
respectively. Passing to inverse limits gives (D4.2), since the
displayed images of compact rank-two modules are closed.

For a single positive prime $p=\ell$, the valuation coordinate is
one and the logarithmic coordinate is zero. If $p\ne\ell$ and
$\ell$ is odd, the logarithm is an isometry on $1+\ell\Z_\ell$, so
\[
 v_\ell(d)=v_\ell(p^{\ell-1}-1)-1.
\]
For $\ell=2$ and $p$ odd, put $\epsilon=(-1)^{(p-1)/2}$.
The logarithm is an isometry on $1+4\Z_2$, and
\[
 v_2(d)=v_2(\epsilon p-1)-2=v_2(p^2-1)-3.
\]
The extra factor two in (D4.2) proves (D4.1). These valuations are
finite because a positive prime is not a root of unity. This proves
openness without discarding the Wieferich cases or the coefficient
prime two. The exact finite-level checks in
\texttt{root\_local\_kummer.py} accompany the general argument.
\end{proof}
